\documentclass[conference]{IEEEtran}
\usepackage[T1]{fontenc}
\usepackage{cite}
\usepackage{amsmath,amssymb}
\usepackage{graphicx}
\usepackage{booktabs}
\usepackage{array}
\usepackage{url}
\usepackage[hidelinks]{hyperref}
\usepackage{xcolor}
\newcommand{\todo}[1]{{\color{red}[#1]}}
\usepackage{float}

% Include image if the file exists, otherwise show a red placeholder box
\newcommand{\figorbox}[2]{%
  \IfFileExists{#2}{\includegraphics[width=#1]{#2}}%
  {\fbox{\parbox{#1}{\centering\vspace{0.6cm}\todo{Missing file: \detokenize{#2}}\vspace{0.6cm}}}}}

\setlength{\textfloatsep}{8pt plus 2pt minus 2pt}
\setlength{\floatsep}{8pt plus 2pt minus 2pt}
\setlength{\intextsep}{8pt plus 2pt minus 2pt}
\begin{document}
\title{Underwater Image Quality Assessment with a Degradation-Aware Multi-Scale Transformer}
\author{\IEEEauthorblockN{Cheemakurthi Swetha, KR Srinidhi Rajamane}
\IEEEauthorblockA{Department of Computer Science and Engineering\\
Shiv Nadar University Chennai, India\\
Registration Number: 24011101016, 24011101059\\
Roll Number: 24110037, 24110033\\
Emails: swetha24110037@snuchennai.edu.in, srinidhirajamane24110033@snuchennai.edu.in}}
\maketitle

\begin{abstract}
Underwater images degrade severely because of turbidity which is due it impurities in the sea water, so a single global score is inefficient because of uneven patches in the images. Our objective is to not only have a  no-reference underwater image quality assessment that predicts human mean opinion scores (MOS) but also show where the quality is lost. We evaluated our model on the UWIQA and SAUD datasets. Our proposed network splits each image into three different scales and encodes them with a shared Swin Transformer. A Transformer provides cross patch context as well as degradation embedding pretrained on haze, color cast and blur conditions using which its tokens through feature wise linear modulation. At the end there are three heads which predict patch the quality, saliency and degradation and fuse their importance through scale-adaptive pooling.Inference combines a layered $K$-fold group , mirrored averaging  and a probability adjustment . On the UWIQA test split the model reached PLCC 0.8265, SROCC 0.8072, KRCC 0.6668 and RMSE 0.0842 , On SAUD the  model reaches PLCC 0.7733, SROCC 0.7499, KRCC 0.5595 and RMSE 11.44 . Our contribution includes a patch weigthed importance and a  degradation conditioned design that produce interpretable quality maps.
\vspace{0.3cm}
\newline\noindent\textbf{Code repository:} \url{https://github.com/Srinidhi-Krishna/UIQA}
\end{abstract}
\vspace{0.3cm}
\begin{IEEEkeywords}
Underwater image quality assessment, no reference, multi-scale Transformer, degradation awareness, UWIQA, SAUD
\end{IEEEkeywords}

\section{Introduction}


Underwater imaging help in conducting marine surveys, subsea inspection and autonomous navigation, but due to wavelength dependent absorption, scattering and turbidity it produces an uneven image. Underwater image quality assessment is hence used to predict human mean opinion scores without a reference. Handcrafted metrics such as UCIQE [2] generalises poorly across water types and global regression networks assume that regions contribute equally to predict quality. Patch-based methods [3] give local estimates but is dependent on generic saliency. transformers capture crosspatch context [5] but they are not designed for underwater degradation.
We implemented a hybrid UIQA pipeline which is evaluated on the UWIQA and SAUD dataset. A swin-t backbone  with squeeze-and-excitation (SE) blocks encodes patches at three scales with shared weights. A transformer with scale embeddings gives context, and a degradation embedding is pretrained on synthetic haze, color cast and blur adjusts its patches through FiLM [9]. three heads estimate patch quality, saliency and degradation, and their combined importance weights guide scale-adaptive pooling. training combines Huber, ranking, PLCC, and degradation losses, flip-based augmentation, and logistic calibration. The model predicts patch level quality, importance  and degradation maps alongside the global score.


\section{Related Works}
\vspace{0.3cm}
\subsection{Underwater Image Quality Assessment}

Because  references are rarely available for underwater hence UIQA is posed as a no reference problem . UCIQE~\cite{yang2015uciqe} combines the standard deviation of chroma, luminance and the saturation. UIQM~\cite{panetta2016uiqm} combines colorfulness sharpness and contrast measures similar to human visual system Both are cheap and explainable but they use fixed manually set weights so they mostly miss localised degradations and generalise poorly across various water types. FDUM~\cite{yang2021uwiqa} combines colorfulness, contrast and sharpness and introduces the UWIQA dataset to be used in this work. NUIQ~\cite{jiang2022nuiq} extracts quality aware features and compresses and regresses them into a score. All these methods score an image which do not emphasize on which regions contribute to the perceived quality.
\vspace{0.3cm}
\subsection{Deep NR-IQA}

Kang et al.~\cite{kang2014} trains CNN on image patches and averages the scores and assumes that every patch contribute equally. WaDIQaM-NR~\cite{bosse2018} learns  weight for every patch and combines  the scores. This is the closest to importance aware pooling but the weights come from generic statistics and the patches are processed independently.TReS~\cite{golestaneh2022tres} includes CNN features with a transformer and also adds relative ranking and self consistency loss to capture global contexts. These models are adapted for natural images and not for the  underwater degradations.
\vspace{0.3cm}
\subsection{EDANet: The Source Paper}

EDANet~\cite{jiang2025eda} is the paper that is most closely aligned to our work. It trains a degradation network DIDNet, on UIEB and predicts the change between a degradation image and its estimated reference one . the decoder fuses into a ResNet50 based network through a degradation aware module and then they are combined globally for scoring. On UWIQA it gives PLCC 0.8156, SROCC 0.8104, KRCC 0.6661, and RMSE 0.0860, and on SAUD PLCC 0.8792, SROCC 0.8517, KRCC 0.7459, and RMSE 9.3911. Its main drawbacks are that the Quality is obtained from globally pooled features and hence there are no explicit per patch importance scores. It makes use of  a single input scale and CNN backbone hence the cross region context is not completely available . Its degradation prior depends on the enhanced references.

\begin{table}[H]
\centering
\caption{Comparison of Method Capabilities}
\label{tab:related}
\footnotesize
\setlength{\tabcolsep}{3pt}
\begin{tabular}{lcccl}
\toprule
Method & Underw. & Local & Context & Cue \\
\midrule
UCIQE/UIQM/FDUM/NUIQ & Yes & No & No & Handcrafted \\
Kang et al. & No & No & No & None \\
WaDIQaM-NR & No  & Yes & No & None \\
RankIQA    & No  & No  & No & None \\
TReS       & No  & No  & Yes & None \\
EDANet     & Yes & Partial & No & Residual \\
\textbf{Ours} & Yes & Yes & Yes & Embedding \\
\bottomrule
\end{tabular}
\end{table}
\vspace{0.3cm}
\subsection{Positioning of This Work}
Table~\ref{tab:related} shows the gap among which the methods are reviewed here while none combines patch level importance weighting and multi scale context and an explicit underwater degradation signal. We keep EDANet's logic that a degradation can help  on synthetic levels. We use it to condition the transformer tokens and also the weight patches at several scale.

\section{Proposed Methodology}
\vspace{0.3cm}
\subsection{Problem Formulation and Dataset}

For an underwater RGB image without original reference the model predicts a global quality score Q[0,1] that approximates to human MOS. It also returns patch level maps for quality qi, importance i, and degradation ri. Only image level MOS is available as human supervised score so each image is a bag of its patches as instances [4]. The patch maps are latent variables learned through backpropagation and not any verified patch level ground truth.
We use two datasets. UWIQA provides 890 underwater images with continuous MOS scores on the [0, 1] scale . Fig.~\ref{fig:classes}(a) shows that the UWIQA labels are concentrated around 0.4 to 0.6 while the classes at 0.8 and above hold only 58 images so a plain sampler would rarely train the model on very good images. Fig.~\ref{fig:mos}(a) confirms the UWIQA label direction that the lowest MOS images are hazy or have a green tint and the high MOS images are clear and sharp. the SAUD is the second benchmark dataset on which EDANet reports its best result. It contains 1000 underwater images with continuous MOS annotations, normalized to [0, 1] for training. Fig.~\ref{fig:classes}(b) shows that SAUD labels are also concentrated around the 0.5, 0.6, and 0.7 classes which have 199, 273, and 201 images respectively  whereas the 0.2 and 1.0 classes hold only 7 and 5 so the extremes range scores are lesser than in UWIQA. Fig.~\ref{fig:mos}(b) confirms the label direction that the lowest MOS images (0.00) have dark tint or hazy while the highest MOS images (0.97–1.00) are clear with natural blue color. Because the SAUD MOS scale differs from UWIQA their predictions are mapped back to the original scale before computing RMSE, so SAUD RMSE is in original MOS units and is not comparable to that of UWIQA RMSE.
The split is made at image level before any patch extraction is being done and hence it is class balanced by MOS class .The same procedure is applied independently to each dataset. With seed 42, a 70/10/20 split is made, and the 20\% part is the fixed test data. The 70\% and 10\% parts of the data is merged  which is repartitioned into K=10 folds. And the following precautions prevent leakage :
No image should appear in more than one split which the code also verifies.
The class balanced  augmentation are applied to the training only.
The degradation estimator is pretrained for each fold on that fold’s training images only.
The logistic calibration is fitted on the out of fold  validation predictions only.
The test split is used to produce the final reported scores.
Near duplicate image pairs are flagged using thumbnail similarity .
SAUD contains the same scene modified by multiple methods so a scene level splitting is applied to ensure that all enhanced versions of the same scene is contain within the same split. This prevents the same scene from appearing in both training as well as the testing thus reducing the risk of data leakage and provide a reliable evaluation of the model's generalization performance.


\begin{figure*}[t]
\centering
\begin{minipage}[t]{0.44\textwidth}
\centering
\includegraphics[width=\linewidth]{Block_dia.png}\\[2pt]
{\footnotesize (a) Model architecture}
\end{minipage}\hfill
\begin{minipage}[t]{0.44\textwidth}
\centering
\includegraphics[width=\linewidth]{Block_dia_2.png}\\[2pt]
{\footnotesize (b) Training, inference, and evaluation pipeline}
\end{minipage}
\caption{ Proposed Model Architecture and Training Pipeline. The model combines multi-scale patches, degradation-aware learning, and Transformer-based quality prediction.}
\label{fig:arch}
\end{figure*}

\begin{figure*}[t]
\centering
\figorbox{0.62\textwidth}{class_distribution.png}\\[1pt]
{\footnotesize (a) UWIQA}\\[4pt]
\figorbox{0.62\textwidth}{class_distribution_saud.png}\\[1pt]
{\footnotesize (b) SAUD}
\caption{MOS Class Distribution. (a) UWIQA and (b) SAUD datasets, showing the original distribution, training distribution, and distribution after class balancing.}
\label{fig:classes}
\end{figure*}
\vspace{0.3cm}
\subsection{Processing Pipeline and Approach}

Fig.~\ref{fig:arch}(a) shows the architecture and Fig.~\ref{fig:arch}(b) the training and inference pipeline.
\vspace{0.2cm}
\subsubsection{Preprocessing and multi scale patches}
Every image is resized to $256\times256$ and normalized with ImageNet statistics. Non-overlapping patches of size $p_s\in\{32,64,128\}$ are extracted (Fig.~\ref{fig:patches}), giving $N_s=(256/p_s)^2$ patches namely 64, 16 and 4 for $N=84$ tokens per image.

\begin{figure*}[t]
\centering
\includegraphics[width=0.55\textwidth]{patch_grids[14].png}
\caption{Multi-scale patch grids on a resized image: $32\times32$ patches ($N=64$), $64\times64$ patches ($N=16$), and $128\times128$ patches ($N=4$).}
\label{fig:patches}
\end{figure*}
\vspace{0.2cm}
\subsubsection{Patch feature extraction}
One ImageNet pretrained ResNet-50~\cite{he2016resnet} is shared across all the scales and are encoded at every patch up to its last convolutional stage. Patches are fed in their original size and are not resized so a $32\times32$ patch gives a $1\times1$ final feature map. an SE block re weights channels and Global average pooling gives a 2048-d vector which a linear layer maps to.
\vspace{0.2cm}
\subsubsection{Degradation estimator and FiLM conditioning}
A small CNN of four stride 2 convolution blocks having group normalization maps the whole image to embeddings $\mathbf{z}\in\mathbb{R}^{128}$ having three levels $\hat{\boldsymbol{\ell}}\in[0,1]^3$ one for haze, color cast  and blur. these three are used because they are controllable and do not need any estimated reference. In stage 1 it is pretrained to reduce the known levels $\boldsymbol{\ell}$ of synthetic degradations drawn uniformly in $[0,1]$  with a mean squared error of
\begin{equation}
\mathcal{L}_{\text{pre}}=\tfrac{1}{3}\lVert \hat{\boldsymbol{\ell}}-\boldsymbol{\ell}\rVert_2^2 .
\end{equation}
Color cast attenuates red the most wheraes haze blends the image towards a blueish green airlight this stage uses no MOS and no estimated reference. In stage 2, $\mathbf{z}$ adjusts every token through FiLM~\cite{perez2018film}:
\begin{equation}
[\boldsymbol{\gamma},\boldsymbol{\beta}]=\mathbf{W}_f\mathbf{z},\qquad
\mathbf{x}_i\leftarrow \mathbf{x}_i\odot(1+\boldsymbol{\gamma})+\boldsymbol{\beta},
\end{equation}
where $\mathbf{W}_f$ is 0 initialized so training the starts from the unconditioned model and  The estimator is then fine tuned with the remaining network. 
\vspace{0.2cm}
\subsubsection{Cross patch and cross scale Transformer}
The input token is $\mathbf{x}_i=\mathbf{f}_i^{(s)}+\mathbf{e}_s+\mathbf{p}_i$, with a learnable embedding $\mathbf{e}_s$ and a position embedding $\mathbf{p}_i$ for each token position within a scale. After FiLM all $N$ tokens pass through prenorm encoder layers with 8 heads and
\begin{equation}
\mathbf{X}'=\mathrm{softmax}\!\left(\frac{\mathbf{Q}\mathbf{K}^{\top}}{\sqrt{d}}\right)\mathbf{V},
\end{equation}
so each patch is provides context to every other patch at every scale. The output is $\mathbf{h}_i\in\mathbb{R}^{d}$.
\vspace{0.2cm}
\subsubsection{Three branch head and importance fusion}
Three MLP branches act on each $\mathbf{h}_i$:
\begin{equation}
q_i=\sigma(g_q(\mathbf{h}_i)),\quad s_i=\sigma(g_s(\mathbf{h}_i)),\quad \mathbf{r}_i=\sigma(g_r(\mathbf{h}_i)),
\end{equation}
where $\mathbf{r}_i\in[0,1]^5$ covers color, haze, blur, contrast and noise. Importance is computed from the token and the branch outputs and are normalized :
\begin{equation}
e_i=\mathrm{MLP}([\mathbf{h}_i;q_i;s_i;d_i]),\qquad
\alpha_i^{(s)}=\frac{\exp e_i}{\sum_{j\in\mathcal{P}_s}\exp e_j}.
\end{equation}
\vspace{0.2cm}
\subsubsection{Scale adaptive pooling}
Each scale is pooled as $Q^{(s)}=\sum_{i\in\mathcal{P}_s}\alpha_i^{(s)}q_i$. The scale weights are combined to learn a global bias and the score is given as :
\begin{equation}
\lambda_s=\mathrm{softmax}_s\!\big(b_s+\mathbf{w}^{\top}\bar{\mathbf{h}}^{(s)}\big),\qquad
\hat{Q}=\sum_{s}\lambda_s Q^{(s)},
\end{equation}
where $\bar{\mathbf{h}}^{(s)}$ is the mean token of scale and the shared  weights is 0 initialized  so that the training starts from equal scale weights.
\vspace{0.2cm}
\subsubsection{Training objective and procedure}
Training minimizes
$\mathcal{L}_{\text{MOS}}$ is a Huber loss ($\delta=0.1$) whose per sample weights grow with the distance from the label  thus extreme quality images are counted more. 
\begin{equation}
\mathcal{L}=\lambda_{1}\mathcal{L}_{\text{MOS}}+\lambda_{2}\mathcal{L}_{\text{rank}}+\lambda_{3}\mathcal{L}_{\text{PLCC}}+\lambda_{4}\mathcal{L}_{\text{deg}}.
\end{equation}$\mathcal{L}_{\text{rank}}=\max(0,m-(\hat{y}_a-\hat{y}_b))$ is a margin loss and $\mathcal{L}_{\text{PLCC}}=1-r(\hat{\mathbf{y}},\mathbf{y})$ is computed per batch and optimizes the correlation calculated. $\mathcal{L}_{\text{deg}}$ is the mean squared error between $\mathbf{r}_i$ and estimated labels. These are not human annotations. The training loader draws MOS classes (Fig.~\ref{fig:classes}) with augmentation thats is flips, 90$^\circ$ rotations The backbone is frozen for the first 20 epochs and then they are fine tuned.
\vspace{0.2cm}
\subsubsection{Inference and calibration}
At the time of testing each model's prediction is averaged over four flip views. A five parameter logistic mapping
is fitted by an unconstrained least squares thereby enabling lowest squared error on the fit data. the fit is not monotonic. We report two scores. The \emph{raw} score is the mean over the $K$ fold models  with no mapping and the \emph{calibrated} score applies the mapping to  model's flip-averaged output and then it averages the $K$ models like in Fig.~\ref{fig:arch}(b). The mapping is fitted separately for each dataset.
\begin{equation}
Q_{\text{fit}}(x)=\beta_1\!\left(\tfrac12-\tfrac{1}{1+e^{\beta_2(x-\beta_3)}}\right)+\beta_4x+\beta_5,
\end{equation}
\vspace{0.2cm}
\subsubsection{Modifications relative to EDANet}
Relative to EDANet (i) the degradation prior is pretrained on synthetic levels instead of pseudo references (ii) it adjusts Transformer tokens through FiLM instead of combining it with the decoder features (iii) quality is combined with the learned patch and its importance is being scaled and (iv) the input is a multiscale transformer which provides context. The components themselves are the existing techniques and our contribution is their incorporation for underwater quality prediction.

\subsection{Implementation Details}

The code uses python, pytorch and torchvision sciPy, and scikit-learn, Table~\ref{tab:hparams} contains the hyperparameter
setting used by both the UIQA and SAUD datasets, which includes the batch size, learning rate, epochs,
and early-stopping. validation and checkpoints use exponential moving average weights, and the best
epoch per fold is selected based on validation SROCC, one iteration represents a complete training run on
a single fold which is continued until the early-stopping patience has reached.

\begin{table}[H]
\centering
\caption{Main hyperparameters.}
\label{tab:hparams}
\footnotesize
\setlength{\tabcolsep}{3pt}
\begin{tabular}{p{0.40\columnwidth}p{0.52\columnwidth}}
\toprule
Item & Value \\
\midrule
Input / patch sizes & $256^2$ / 32, 64, 128 (native size) \\
Embedding dim. $d$ / FiLM dim. & 256 / 128 \\
Transformer & 4 layers, 8 heads, FF 512, dropout 0.4 \\
Optimizer & AdamW, lr $3\!\times\!10^{-5}$, weight decay $10^{-3}$ \\
Schedule & Cosine with warm restarts ($T_0{=}50$, $T_{\text{mult}}{=}2$, floor $3\!\times\!10^{-7}$) \\
Backbone & frozen 20 epochs, then lr $\times0.2$ \\
Batch size / batches per epoch & 16 / 25 \\
Max epochs / early stop & 200 / patience 30 on val SROCC, counted from backbone unfreezing \\
EMA decay / gradient clip & 0.995 / 1.0 \\
Huber $\delta$ / margin $m$ / $\gamma$ & 0.1 / 0.05 / 0.5 \\
Loss weights $\lambda_{1..4}$ & 1.0, 0.5, 0.5, 0.1 \\
Estimator pretraining & 8 epochs $\times$ 100 iters, lr $3\!\times\!10^{-4}$ \\
Test split / folds $K$ & 20\% held out / 10 on the 80\% pool \\
Seed & 42 (split and folds fixed) \\
\bottomrule
\end{tabular}
\end{table}

\section{Results and Discussion}
\vspace{0.3cm}
\subsection{Experimental Setup and Evaluation}

We measure 4 metrics between predictions $\hat{y}$ and MOS $y$. PLCC is used for linear correlation, SROCC and KRCC measures rank correlation and RMSE measures absolute error on the MOS scale for  each dataset. validation data is restricted to early stopping and calibration. final scores are on the  20\% test split of each dataset (179 UWIQA images and 210 SAUD images) using the $K$-fold cross validation and we measure both the raw and calibrated scores.The baseline values are the ones published for EDANet.
\vspace{0.3cm}
\subsection{Quantitative and Qualitative Results}

\begin{table}[H]
\centering
\caption{Performance comparison across Datasets }
\label{tab:results}
\footnotesize
\setlength{\tabcolsep}{3pt}
\begin{tabular}{>{\raggedright\arraybackslash}p{0.34\columnwidth}cccc}
\toprule
Method & PLCC$\uparrow$ & SROCC$\uparrow$ & KRCC$\uparrow$ & RMSE$\downarrow$ \\
\midrule
\multicolumn{5}{l}{\emph{UWIQA}} \\
EDANet~\cite{jiang2025eda} & 0.8156 & \textbf{0.8104} & 0.6661 & 0.0860 \\
Ours (raw) & 0.8248 & 0.8072 & 0.6666 & 0.0867 \\
Ours (calibrated) & \textbf{0.8265} & 0.8072 & \textbf{0.6668} & \textbf{0.0842} \\
$\Delta$ vs.\ EDANet & +0.0109 & $-$0.0032 & +0.0007 & $-$0.0018 \\
\midrule
\multicolumn{5}{l}{\emph{SAUD}} \\
EDANet~\cite{jiang2025eda} & \textbf{0.8792} & \textbf{0.8517} & \textbf{0.7459} & \textbf{9.3911} \\
Ours (raw) & 0.7638 & 0.7498 & 0.5594 & 11.62$^\dagger$ \\
Ours (calibrated) & 0.7733 & 0.7499 & 0.5595 & 11.44$^\dagger$ \\
$\Delta$ vs.\ EDANet & $-$0.1059 & $-$0.1018 & $-$0.1864 & +2.05$^\dagger$ \\
\bottomrule
\end{tabular}
\end{table}

Table~\ref{tab:results} gives the main comparison of the metric values obtained. On UWIQA our model has PLCC of 0.8265, KRCC of
0.6668, and RMSEof 0.0842 , while EDANet has a higher value on SROCC which is 0.8104. This
suggests that our model performs relatively better than EDANet for UWIQA. On SAUD our model reaches PLCC 0.7733, SROCC 0.7499, KRCC 0.5595, and RMSE
11.44 whereas EDANet’s published values (0.8792, 0.8517, 0.7459, and 9.3911) are better on 
metrics. The gaps are 0.1059 in PLCC, 0.1018 in SROCC, and 0.1864 in KRCC. We therefore do not claim any advantage on SAUD


\begin{figure}[H]
\centering
\figorbox{0.95\columnwidth}{mos_direction_check.png}\\[1pt]
{\footnotesize (a) UWIQA}\\[4pt]
\figorbox{0.95\columnwidth}{mos_direction_check_saud.png}\\[1pt]
{\footnotesize (b) SAUD}
\caption{MOS Quality Comparison. (a) UWIQA and (b) SAUD, showing low- and high-MOS images.
}
\label{fig:mos}
\end{figure}
\vspace{0.3cm}
\subsection{Discussion and Limitations}

Calibration mainly reduces prediction error and improves the similarity between predicted scores and human MOS. On UWIQA, calibration improves PLCC and reduces RMSE, while SROCC and KRCC is almost the same. This is because adjustment is done to predicted scores to match the MOS range without drastically changing their ranking. even before calibration, the model performs better than EDANet in PLCC and KRCC, showing that these improvements are not because of calibration. On SAUD, calibration produces similar improvements, but the model still performs worse than EDANet, telling that calibration does not resolve the performance gap.

Compared with EDANet on UWIQA, the proposed model achieves higher PLCC and lower RMSE, indicating better agreement with human MOS values. However, the slightly lower SROCC suggests that image ranking still needs improvement. since these differences are small and negligible, the model is comparable with EDANet on UWIQA on the basis of performance.on SAUD, the predictions are concentrate around the middle of the quality scale, causing for both high quality and poor-quality  images to receive poorly predicted scores.the largest prediction errors occur mainly at these extremes.This suggests that the model may have difficulty distinguishing between modified versions of the same scene, particularly when their quality differences are mainly due to color rendering. However, further experiments are required to confirm this observation.
\vspace{0.3cm}

The following limitations apply regardless of the final scores.
\vspace{0.3cm}
\begin{itemize}
\item There are no patch level MOS labels, so the quality, importance and degradation maps are not physical measurements of absorption or turbidity.
\item The degradation branch is supervised by estimated labels, and the estimator uses a synthetic pretraining task that is different from EDANet, which are not yet validated.
\item TThe backbone is pretrained on imageNet images, which have different characteristics from underwater image.
\item A 32×32 patch gives a 1×1 ResNet 50 feature map, so it learns little spatial information.
\item UWIQA and SAUD labels concentrated in the mid range Fig.~\ref{fig:classes}(Fig. 2), and the top and bottom few classes have few images , so scores for the extreme class range images may reduce accuracy. The SAUD predictions show this directly through their range compression
\item The reported scores come from one training seed per dataset, so we cannot separate the small differences from EDANet on UWIQA from variation across training seeds
\item SAUD and UWIQA use different MOS scales, so we cant compare the RMSE metric across datasets, and the EDANet protocol on SAUD may differ from ours
\item The model is clearly below EDANet’s published results on SAUD, so the claims of this pipeline is supported on UWIQA only.
\end{itemize}
\vspace{0.2cm}
\section{Conclusion and Future Work}

Our project implements a degradation-aware, multi-scale transformer for underwater image quality
assessment that learns patch quality and importance from image level MOS. On UWIQA the model
performance reaches PLCC 0.8265, SROCC 0.8072, KRCC 0.6668, and RMSE 0.0842, which is higher in
PLCC and lower in RMSE than the published EDANet values and slightly lower in SROCC. On SAUD it
reaches PLCC 0.7733, SROCC 0.7499, KRCC 0.5595, and RMSE 11.44 which is below EDANet’s
published 0.8792, 0.8517, 0.7459, and 9.3911 on every metric, so the design does not yet transfer to SAUD
this is mainly because the maps are latent attributions and the degradation supervision is synthetic or
handcrafted and also the the enhancements of the SAUD dataset could not be able to completely
implement into the model
conclusions about where quality is lost need human validation.The future steps that can be taken to
improve are:
(i) running the ablations over multiple seeds to estimate variance
(ii) reducing the range compression seen on SAUD
(iii) improving ranking consistency, including overlapping patches to reduce boundary effects.
\vspace{0.3cm}
\begin{thebibliography}{10}

\bibitem{yang2015uciqe}
M. Yang and A. Sowmya, ``An underwater color image quality evaluation metric,''
\emph{IEEE Trans. Image Process.}, vol. 24, no. 12, pp. 6062--6071, 2015.

\bibitem{bosse2018}
S. Bosse, D. Maniry, K.-R. M\"uller, T. Wiegand, and W. Samek, ``Deep neural networks for no-reference and full-reference image quality assessment,''
\emph{IEEE Trans. Image Process.}, vol. 27, no. 1, pp. 206--219, 2018.

\bibitem{ilse2018}
M. Ilse, J. Tomczak, and M. Welling, ``Attention-based deep multiple instance learning,''
in \emph{Proc. Int. Conf. Mach. Learn. (ICML)}, 2018, pp. 2127--2136.

\bibitem{dosovitskiy2021}
A. Dosovitskiy \emph{et al.}, ``An image is worth 16$\times$16 words: Transformers for image recognition at scale,''
in \emph{Proc. Int. Conf. Learn. Represent. (ICLR)}, 2021.

\bibitem{yang2021uwiqa}
N. Yang, Q. Zhong, K. Li, R. Cong, Y. Zhao, and S. Kwong, ``A reference-free underwater image quality assessment metric in frequency domain,''
\emph{Signal Process.: Image Commun.}, vol. 94, Art. no. 116218, 2021.

\bibitem{he2016resnet}
K. He, X. Zhang, S. Ren, and J. Sun, ``Deep residual learning for image recognition,''
in \emph{Proc. IEEE Conf. Comput. Vis. Pattern Recognit. (CVPR)}, 2016, pp. 770--778.

\bibitem{perez2018film}
E. Perez, F. Strub, H. de Vries, V. Dumoulin, and A. Courville, ``FiLM: Visual reasoning with a general conditioning layer,''
in \emph{Proc. AAAI Conf. Artif. Intell.}, 2018.

\bibitem{jiang2025eda}
Q. Jiang, Y. Gu, Z. Wu, C. Li, H. Xiong, F. Shao, and Z. Wang, ``Deep underwater image quality assessment with explicit degradation awareness embedding,''
\emph{IEEE Trans. Image Process.}, vol. 34, pp. 1297--1311, 2025.

\bibitem{panetta2016uiqm}
K. Panetta, C. Gao, and S. Agaian, ``Human-visual-system-inspired underwater image quality measures,''
\emph{IEEE J. Ocean. Eng.}, vol. 41, no. 3, pp. 541--551, 2016.

\bibitem{jiang2022nuiq}
Q. Jiang, Y. Gu, C. Li, R. Cong, and F. Shao, ``Underwater image enhancement quality evaluation: Benchmark dataset and objective metric,''
\emph{IEEE Trans. Circuits Syst. Video Technol.}, vol. 32, no. 9, pp. 5959--5974, 2022.

\bibitem{kang2014}
L. Kang, P. Ye, Y. Li, and D. Doermann, ``Convolutional neural networks for no-reference image quality assessment,''
in \emph{Proc. IEEE Conf. Comput. Vis. Pattern Recognit. (CVPR)}, 2014, pp. 1733--1740.

\bibitem{golestaneh2022tres}
S. A. Golestaneh, S. Dadsetan, and K. M. Kitani, ``No-reference image quality assessment via transformers, relative ranking, and self-consistency,''
in \emph{Proc. IEEE/CVF Winter Conf. Appl. Comput. Vis. (WACV)}, 2022, pp. 1220--1230.

\end{thebibliography}

\end{document}
